% !TeX root = ../凸优化笔记.tex
% 由主文档通过 \input 载入；请编译主文档。
\chapter*{课程大纲与阅读说明}
\addcontentsline{toc}{chapter}{课程大纲与阅读说明}

\section*{课程定位}

凸优化（convex optimization）研究在凸可行集上极小化凸函数的问题。本书面向学过线性代数与多元微积分的本科生，围绕三个问题展开：怎样建立凸模型，怎样证明候选解最优，以及怎样以有限精度求解。基本概念和核心算法给出推导；涉及较多凸分析工具的进阶定理注明引用或证明思路，供第二轮学习使用。

全书沿用 Boyd 和 Vandenberghe 的《Convex Optimization》\cite{boyd-vandenberghe} 的理论、应用、算法主线，以最小二乘、资源分配和几何投影贯穿各章。概率知识主要用于统计应用，数值线性代数按算法需要补充。

\section*{全书课程大纲}
\begin{enumerate}
  \item 第\ref{ch:co-introduction}章：绪论与优化建模。变量、目标、约束、可行性、最优值与数值精度。
  \item 第\ref{ch:co-sets}章：凸集。凸组合、凸包、典型凸集、保凸运算、投影、分离与对偶锥。
  \item 第\ref{ch:co-functions}章：凸函数。凸性判据、保凸运算、共轭、次梯度与拟凸性。
  \item 第\ref{ch:co-models}章：凸优化问题与标准模型。最优性、等价变换、线性规划、二次规划及锥规划。
  \item 第\ref{ch:co-duality}章：拉格朗日对偶与最优性条件。强弱对偶、Slater 条件、KKT 条件与灵敏度。
  \item 第\ref{ch:co-fitting}章：逼近、拟合与正则化。误差度量、稀疏拟合、鲁棒约束与形状限制。
  \item 第\ref{ch:co-statistics}章：统计估计与实验设计。似然、后验、熵、检测器和信息矩阵。
  \item 第\ref{ch:co-geometry}章：几何优化与分类。投影、距离、椭球、间隔与定位。
  \item 第\ref{ch:co-unconstrained}章：梯度方法与牛顿方法。线搜索、梯度法、投影与近端方法、次梯度法及牛顿法。
  \item 第\ref{ch:co-equality}章：等式约束优化算法。消元、约束牛顿法、残差与块线性系统。
  \item 第\ref{ch:co-interior}章：内点法。障碍函数、中心路径、可行性与原对偶迭代。
\end{enumerate}
附录\ref{app:co-background}补充数学基础，附录\ref{app:co-quadratic}讨论两个二次函数相关问题，附录\ref{app:co-numerical}说明数值线性代数。

\section*{阅读与练习}

首先顺序学习第\ref{ch:co-introduction}至\ref{ch:co-duality}章。应用部分的三个章节可以按兴趣选读；算法部分从梯度法及其处理简单约束、非光滑项的延伸出发，再学习无约束、等式约束和不等式约束牛顿方法。相对内部、矩阵凸性、自协调分析和附录\ref{app:co-quadratic}属于进阶内容，不影响首次掌握主要方法。

阅读时须分别回答：结论的假设是什么，关键不等式来自哪里，去掉某个假设是否仍然成立。各章习题给出必要提示或核对结果，完整推导仍应独立完成。数值练习要求同时检查目标值、可行性和适用的最优性指标；实验结果不能替代理论证明。

本书所说的“理解一个方法”，包括能用自己的话解释它的目标，能在一个两变量问题上手算，并能指出它何时不适用。例如，学完凸集应能判断两种可行方案的平均是否仍可行；学完凸函数应能解释为什么切线给出下界；学完对偶应能给一个候选解附上数值相等的下界；学完算法应能说明停止条件究竟控制什么误差。公式是这些判断的简明表达，而不是需要孤立记忆的规则。

首次学习可将每章的定义、基本定理、数值例题和习题作为主线。遇到一般分离定理、共轭双重化、自协调性等选读内容，先理解结论的用途及假设，再回到主线；不必为了尚未用到的推广中断学习。附录中的矩阵知识也可按需查阅。例如第一次看到 $P\succeq0$ 时，只需先记住它等价于所有方向上的二次型 $d^{\mathsf T}Pd$ 非负。

